\SourceLeaf{351}
In meiner (der ersten) Construction, wo $\delta$ unbestimmt, ist die Clise wie
\[
(a-b)(a-c)(b-c).
\]
Nennt man $e$ den vierten Schnitt[punkt] der Ellipse mit dem durch $b$, $c$, $d$ gelegten Kreise und setzt
\[
e = \varepsilon m + \frac{n}{\varepsilon},
\]
wodurch
\[
\beta\gamma\delta\varepsilon = 1
\]
wird, so ist, in \textsc{Newton}s Construction, die Clise der Differenz zweier Werthe von $t$ wie
\[
(c-b)(b-e)(c-e).
\]
\Needspace{10\baselineskip}Auch kann obige Formel ($a=\varepsilon\rho$) besser so geschrieben werden:
\[
t = -\frac{(m-\delta\beta n)(m-\delta\gamma n)(m-\delta\varepsilon n)\rho}{\delta(1-\rho)mn} + \frac{(\beta+\gamma)mm-\dfrac{1}{\varepsilon}mn+\delta nn}{m}
\]
\[
= -\frac{(b-c)(b-e)(c-e)}{\delta(\beta-\gamma)(\beta-\varepsilon)(\gamma-\varepsilon)}\cdot\frac{\rho}{(1-\rho)mn} + \frac{(\beta+\gamma)mm-\dfrac{1}{\varepsilon}mn+\delta nn}{m}.
\]

\Needspace{10\baselineskip}\begin{center}\textbf{BEMERKUNGEN.}\end{center}

Die Notiz [III], die auf einem einzelnen Blatte steht, gibt den Inhalt von Untersuchungen wieder, über die \textsc{Gauss} in dem S.\,295 bis 296 dieses Bandes abgedruckten Briefe an \textsc{Schumacher} vom 12. Mai 1843 berichtet hat; auf diesen möge zur Erläuterung des Textes verwiesen werden. Die Notizen [I] und [II], die mit [III] in engem Zusammenhange stehen, finden sich in einem Handbuche und auf einem Zettel; [I] stammt wahrscheinlich bereits aus dem Jahre 1831, [II] wohl aus der Zeit um 1840.
\hfill \textsc{Stäckel}.
